% ===================================================================
%  TABELUL Nr. 5 — Casa și cum se construiește.
%  Traduction 2026 d'apres texto/io, controlee sur texto/fr, texto/en
%  et texto/es. Consigne : voir texto/ro/10-tabelul-01.tex.
%
%  Le tableau 5 est un DIALOGUE, et les deux volumes composent le nom
%  du parleur en petites capitales : on garde \textsc{}.
%
%  Les renvois du plan sont des LETTRES — (a) le couloir, (b) le
%  salon... — et non des numeros, y compris la ou l'imprimeur a
%  compose « (1) » pour « (l) ».
%
%  LE TABLEAU DES METIERS, ET C'EST ICI QUE LE PARTAGE DU LEXIQUE SE
%  VOIT LE MIEUX. Le macon est « zidar », de « zid » le mur, qui est
%  slave ; le charpentier « dulgher », qui est hongrois ; le forgeron
%  « fierar », de « fier », qui est latin ; le menuisier « tâmplar »,
%  qui est hongrois aussi. Quatre corps de metier sur le meme
%  chantier, trois origines.
% ===================================================================

%%K t05-apar-1 apar
\VUsaut{6.0mm}
\VUtitre{15.0pt}{60}{TABELUL Nr. 5}
\VUsaut{1.4mm}
\VUtitre{11.4pt}{0}{Casa și cum se construiește.}
\VUsaut{2.2mm}

%%K t05-tit-1 sub
\VUpk{10.2pt}{40}{Întâlnire fericită.}
\VUsaut{3.0mm}

%%K t05-01-1 p
\textsc{Ioan}. --- Unchiule, aș vrea tare mult să aflu cum se construiește o
\VUgras{casă}.

%%K t05-01-2 p
\textsc{Unchiul} \textsuperscript{(80)}. --- Nu este greu, băiete; vino cu mine și
îți voi arăta pe cea pe care o construiește domnul Bernard \textsuperscript{(79)}
și în care voi locui.

%%K t05-01-3 p
\textsc{Ioan}. --- Și cine a desenat \VUgras{planul} \textsuperscript{(76)} acestei
case?

%%K t05-01-4 p
\textsc{Unchiul}. --- \VUgras{Arhitectul} \textsuperscript{(75)}; a făcut un desen
foarte limpede, a fost încuviințat, iar \VUgras{antreprenorul}
\textsuperscript{(77)} a început lucrările.

%%K t05-01-5 p
\textsc{Ioan}. --- Și cine este antreprenorul?

%%K t05-01-6 p
\textsc{Unchiul}. --- Domnul Blanc, tatăl tovarășului tău Iacob. El conduce
lucrătorii împreună cu domnul Roux, arhitectul.

%%K t05-01-7 p
Iată-l! Domnul Blanc vine tocmai la vreme cu \VUgras{rigla} sa
\textsuperscript{(78)}, iar domnul Roux cu planul său.

%%K t05-01-8 p
Bună ziua, domnilor. Ce mai faceți?

%%K t05-01-9 p
\textsc{Domnul Roux}. --- Foarte bine, mulțumesc. Bună ziua, Ioane; cum a mai
crescut băiatul!

%%K t05-01-10 p
\textsc{Unchiul}. --- Într-adevăr, un om mic de tot. Este foarte serios; trebuie
să i se lămurească tot ce vede. Astăzi vrea să afle cum se construiește o casă.

%%K t05-tit-2 sub
\VUpk{10.2pt}{40}{Lucrătorii.}
\VUsaut{3.0mm}

%%K t05-01-11 p
\textsc{Domnul Blanc}. --- Vino atunci cu mine, tinere prieten, și îți voi lămuri
îndată totul, căci îmi plac mult copiii care vor să învețe.

%%K t05-01-12 p
Vezi lucrătorul acela cu \VUgras{lopata} \textsuperscript{(82)} în mâini: este
\VUgras{săpătorul} \textsuperscript{(81)}. Sapă \VUgras{șanțul}
\textsuperscript{(46)}, ca să pună țeava de apă. Tot el a scos pământul acolo unde
stă casa, ca zidarul să poată ridica cei patru \VUgras{pereți}. Partea acestor
pereți care este în pământ se numește \VUgras{temelie} \textsuperscript{(2)}.
Spațiul închis de temelie alcătuiește o \VUgras{pivniță} \textsuperscript{(3)}.
Această pivniță are o ferestruică, numită \VUgras{răsuflătoare}
\textsuperscript{(5)}, o \VUgras{ușă} \textsuperscript{(4)} și o scară.

%%K t05-01-13 p
\VUgras{Zidarul} \textsuperscript{(108)} a urmat să ridice pereții, lăsând goluri
pentru \VUgras{uși} \textsuperscript{(8)} și pentru \VUgras{ferestre}
\textsuperscript{(17)}.

%%K t05-01-14 p
\textsc{Ioan}. --- Dar nu îl văd, pe zidarul acesta!

%%K t05-01-15 p
\textsc{Domnul Blanc}. --- A sfârșit deja de lucrat la casă. Este colo, lângă
\VUgras{grajd} \textsuperscript{(54)}, pe \VUgras{schela} sa
\textsuperscript{(110)}, și ridică peretele \VUgras{spălătoriei}
\textsuperscript{(56)}.

%%K t05-01-16 p
Are o mistrie, o troacă și un \VUgras{fir cu plumb} \textsuperscript{(109)}. Dar
numele acestea nu le vei ține minte.

%%K t05-01-17 p
\textsc{Ioan}. --- Ba da, le voi ține minte, fiindcă îi voi cere unchiului o
mistrie mică și o troacă, iar firul cu plumb îl voi face singur dintr-o sforicică.

%%K t05-tit-3 sub
\VUsaut{3.0mm}
\VUpk{10.2pt}{40}{Materialele.}
\VUsaut{3.0mm}

%%K t05-01-18 p
\textsc{Unchiul}. --- Bine. Dar spune-mi, Ioane, ce trebuie ca să construiești o
casă?

%%K t05-01-19 p
\textsc{Ioan}. --- Trebuie \VUgras{piatră cioplită} \textsuperscript{(59)} și
\VUgras{mortar} \textsuperscript{(65)}.

%%K t05-01-20 p
\textsc{Unchiul}. --- Chiar așa. Privește omul acela cu pălărie mare, pe
\VUgras{căruța} lui \textsuperscript{(136)}. Este \VUgras{cărăușul}
\textsuperscript{(135)}. În fiecare zi aduce aici \VUgras{bolovani de piatră}
\textsuperscript{(60)}. Descarcă în curte, iar \VUgras{pietrarii}
\textsuperscript{(83)} cioplesc bolovanii și îi taie cu \VUgras{fierăstrăul lor
de piatră} \textsuperscript{(85)}. \VUgras{Lucrătorul ajutor}
\textsuperscript{(146)} îi duce zidarului, iar acesta îi pune la loc.

%%K t05-01-21 p
Între timp \VUgras{amestecătorul} \textsuperscript{(87)} pregătește mortarul. Îi
trebuie \VUgras{nisip} \textsuperscript{(62)}, \VUgras{var}
\textsuperscript{(63)} și \VUgras{apă} \textsuperscript{(64)}. Varul stă lângă el
într-un \VUgras{sac} \textsuperscript{(71)}; \VUgras{ucenicul zidarului}
\textsuperscript{(111)} îi aduce nisip în \VUgras{roabă} \textsuperscript{(112)},
iar \VUgras{ucenicul} \textsuperscript{(113)} stă la pompă
\textsuperscript{(58)}. Umple \VUgras{găleata} \textsuperscript{(114)}. Mortarul va
fi gata îndată. \VUgras{Cel care dă} \textsuperscript{(107)} îl ia și îl duce pe
scară la schelă, unde zidarul isprăvește această latură a casei.

%%K t05-01-22 p
Și acum: cum se numește lucrătorul care face \VUgras{acoperișul}
\textsuperscript{(31)}?

%%K t05-01-23 p
\textsc{Ioan}. --- Este \VUgras{învelitorul} \textsuperscript{(118)}.

%%K t05-tit-4 sub
\VUsaut{3.0mm}
\VUpk{10.2pt}{40}{Acoperișul casei.}
\VUsaut{3.0mm}

%%K t05-01-24 p
\textsc{Domnul Blanc}. --- Da, dar mai întâi vine \VUgras{dulgherul}
\textsuperscript{(115)}.

%%K t05-01-25 p
Dulgherul face \VUgras{căpriorii} \textsuperscript{(35)}. Leagă \VUgras{grinzile}
\textsuperscript{(37)} cu \VUgras{cepuri} \textsuperscript{(117)}. Lucrătorii
aceia de sus, în pantaloni largi de catifea, sunt dulgheri. Unul ține o bucată de
lemn, celălalt cioplește un cep cu \VUgras{securea} sa \textsuperscript{(116)}. Cel
de lângă \VUgras{coș} \textsuperscript{(41)} este învelitorul. Bate șipcile de
\VUgras{grinzi} (*), iar \VUgras{ardezia} \textsuperscript{(66)} de șipci. Uneori
pune țiglă.

%%K t05-noto-1 noto
\VUnotes{143.2mm}{%
(*) \VUgras{Balk-o} = germ. \textit{Balken}, engl. \textit{joist}, fr.
\textit{solive}, ital. \textit{travicello}, rom. \textit{grindă}, sp. și port.
\textit{viga}. Rădăcină tehnică, încă neprimită, dar propusă aici pentru a reda
înțelesul cuvântului francez \textit{solive}, care nu este nici \textit{trabo}
(fr. \textit{poutre}), nici bucată de lemn. Este un buștean cioplit care poartă
podeaua și tavanul.%
}

%%K t05-01-26 p
Acoperișul grajdului este de \VUgras{țiglă} \textsuperscript{(55)}, cel al casei
de \VUgras{ardezie} \textsuperscript{(32)}, iar cel al verandei de \VUgras{zinc}
\textsuperscript{(24)}.

%%K t05-01-27 p
\textsc{Ioan}. --- Acoperișurile de zinc le face tot învelitorul?

%%K t05-01-28 p
\textsc{Domnul Blanc}. --- Nu, este \VUgras{tinichigiul} \textsuperscript{(121)}
sau \VUgras{instalatorul}, acela care pune \VUgras{lucarna}
\textsuperscript{(38)} în acoperișul casei. Tot el pune \VUgras{jgheaburile}
\textsuperscript{(43)} și \VUgras{burlanele} \textsuperscript{(42)}. Lipește
zincul și tabla. El a prins \VUgras{paratrăsnetul} \textsuperscript{(40)} și
\VUgras{girueta} \textsuperscript{(39)} pe \VUgras{fronton}
\textsuperscript{(36)}.

%%K t05-01-29 p
\textsc{Ioan}. --- Așadar casa este gata?

%%K t05-tit-5 sub
\VUsaut{3.0mm}
\VUpk{10.2pt}{40}{Uneltele.}
\VUsaut{3.0mm}

%%K t05-01-30 p
\textsc{Domnul Blanc}. --- Nu tocmai. \VUgras{Ipsosarul} \textsuperscript{(99)}
ridică din \VUgras{cărămidă} \textsuperscript{(61)} pereții despărțitori, care o
împart în odăi felurite. Pune \VUgras{tencuiala} \textsuperscript{(70)} pe
\VUgras{tavane} \textsuperscript{(26)} și pe pereți. Acum lucrează la tavanul
\VUgras{verandei} \textsuperscript{(33)}. Ca și zidarul, are o \VUgras{mistrie}
\textsuperscript{(100)} și o \VUgras{troacă} \textsuperscript{(101)}.

%%K t05-01-31 p
Apoi tâmplarul face ușile, ferestrele, \VUgras{dușumelele}
\textsuperscript{(16)} și \VUgras{obloanele} \textsuperscript{(19)}.

%%K t05-01-32 p
Acum lucrează în curte la \VUgras{tejgheaua} sa \textsuperscript{(90)}. Are un
\VUgras{fierăstrău} \textsuperscript{(94)}, ca să taie \VUgras{lemnul}
\textsuperscript{(69)}, o \VUgras{rindea} \textsuperscript{(91)}, o
\VUgras{menghină} \textsuperscript{(92)}, o \VUgras{daltă} \textsuperscript{(94
\textit{bis})}, un \VUgras{compas} \textsuperscript{(95)} și un \VUgras{ciocan de
lemn} \textsuperscript{(95 \textit{bis})}.

%%K t05-01-33 p
\textsc{Ioan}. --- Ah, să fii tâmplar este și mai vesel decât să fii zidar!
Unchiule, îmi cumperi o tejghea, un fierăstrău și o rindea? Voi face o cușcă
pentru Grivei.

%%K t05-01-34 p
\textsc{Unchiul}. --- Ți le cumpăr, băiete.

%%K t05-01-35 p
\textsc{Ioan}. --- Ce bucuros sunt! Și cine este omul acela care stă la ușa de
intrare?

%%K t05-tit-6 sub
\VUsaut{3.0mm}
\VUpk{10.2pt}{40}{Lucrările de zugrăveală.}
\VUsaut{3.0mm}

%%K t05-01-36 p
\textsc{Domnul Blanc}. --- Este \VUgras{zugravul} \textsuperscript{(102)}. Stă pe
scara lui cu o oală de \VUgras{vopsea} \textsuperscript{(103)} și cu
\VUgras{bidineaua} sa \textsuperscript{(104)}. Acoperă lemnul ușii de intrare, ca
să țină mai mult.

%%K t05-01-37 p
Tot el lipește tapetele în odăi.

%%K t05-01-38 p
\VUgras{Tapițerul} \textsuperscript{(107)} pe \VUgras{scărița}
\textsuperscript{(106)}, pe \VUgras{balcon} \textsuperscript{(28)}, atârnă
\VUgras{storul} \textsuperscript{(29)}.

%%K t05-01-39 p
Lucrătorul care stă la fereastra cea mare este \VUgras{geamgiul}
\textsuperscript{(96)}. Pune \VUgras{geamurile} \textsuperscript{(18)} în
\VUgras{chit} \textsuperscript{(73)}. Iar lucrătorul acela care stă în genunchi la
ușa pivniței este \VUgras{lăcătușul} \textsuperscript{(97)}. Pune
\VUgras{broasca} \textsuperscript{(6)}. \VUgras{Pila} \textsuperscript{(98)} stă
lângă el.

%%K t05-01-40 p
\textsc{Ioan}. --- Iar acela care bate cu \VUgras{ciocanul}
\textsuperscript{(84)} într-o bucată de fier este tot lăcătuș?

%%K t05-01-41 p
\textsc{Domnul Blanc}. --- Mai bine zis, este \VUgras{fierar}
\textsuperscript{(125)}. Făurește \VUgras{piesele de fier} \textsuperscript{(67)}
pentru \VUgras{electrician} \textsuperscript{(124)}, care stă pe acoperișul
\VUgras{remizei} \textsuperscript{(51)}. Are o \VUgras{vatră}
\textsuperscript{(126)}, o \VUgras{nicovală} \textsuperscript{(127)} și
\VUgras{clești} \textsuperscript{(128)}.

%%K t05-01-42 p
Electricianul prinde \VUgras{sârma de aramă} \textsuperscript{(44)} a telefonului.

%%K t05-tit-7 sub
\VUpk{10.2pt}{40}{Lucrul în grădină.}
\VUsaut{3.0mm}

%%K t05-01-43 p
\textsc{Unchiul}. --- În fundul \VUgras{curții} \textsuperscript{(45)}, lângă
\VUgras{gratii} \textsuperscript{(47)} și \VUgras{poartă} \textsuperscript{(48)},
îl vezi pe \VUgras{grădinar} \textsuperscript{(129)}. Pregătește o \VUgras{grădină}
mică \textsuperscript{(49)}. A săpat pământul cu \VUgras{cazmaua} sa
\textsuperscript{(131)}, seamănă semințe și netezește cu \VUgras{grebla}
\textsuperscript{(130)}. Apoi din \VUgras{stropitoarea} sa \textsuperscript{(132)}
va uda \VUgras{straturile} \textsuperscript{(150)}, și plantele vor crește.

%%K t05-01-44 p
\VUgras{Îngrijitorul de cai} \textsuperscript{(133)}, care tocmai a țesălat calul
și i-a dat nutreț, curăță \VUgras{trăsura} \textsuperscript{(52)} cu buretele, cu
\VUgras{pompa de mână} \textsuperscript{(134)} și cu o găleată de apă. Apoi o va
duce în remiză și va încuia ușa cu un \VUgras{zăvor} greu \textsuperscript{(53)}.

%%K t05-01-45 p
Iată, acum ești mulțumit, cred; lămuriri au fost destule.

%%K t05-01-46 p
\textsc{Ioan}. --- Da, unchiule, dar aș vrea să văd casa pe dinăuntru.

%%K t05-01-47 p
\textsc{Unchiul}. --- Aceasta va fi mâine. Domnul Roux îți va lămuri planul și va
numi fiecare odaie.

%%K t05-tit-8 sub
\VUsaut{3.0mm}
\VUpk{10.2pt}{40}{Împărțirea casei.}
\VUsaut{2.4mm}

%%K t05-apar-2 apar
{\centering\textit{(Vezi planul.)}\par}
\VUsaut{2.4mm}

%%K t05-01-48 p
\textsc{Arhitectul}. --- Să mergem, Ioane, te voi purta prin părțile felurite ale
casei.

%%K t05-01-49 p
Să intrăm la \VUgras{parter} \textsuperscript{(7)}. Iată butonul
\VUgras{soneriei} \textsuperscript{(11)}. Urcăm cele trei \VUgras{trepte}
\textsuperscript{(14)} ale \VUgras{peronului} \textsuperscript{(9)}, trecem
\VUgras{pragul} \textsuperscript{(10)} \VUgras{ușii de intrare}
\textsuperscript{(8)} --- și iată-ne la piciorul \VUgras{scării}
\textsuperscript{(13)}.

%%K t05-01-50 p
\textsc{Ioan}. --- Acesta este coridorul.

%%K t05-01-51 p
\textsc{Arhitectul}. --- Nu, acesta este \VUgras{vestibulul}
\textsuperscript{(12)}. \VUgras{Coridorul} \textsuperscript{(a)} este în fund. La
dreapta --- \VUgras{salonul} \textsuperscript{(b)} și \VUgras{sufrageria}
\textsuperscript{(c)}; în fața sufrageriei --- \VUgras{bucătăria}
\textsuperscript{(d)} și \VUgras{oficiul} \textsuperscript{(e)}; apoi, dinspre
fațadă, \VUgras{biroul} \textsuperscript{(f)}. \VUgras{Veranda}
\textsuperscript{(23)} cu \VUgras{ușa} ei \VUgras{cu geam}
\textsuperscript{(25)} se lipește de salon.

%%K t05-01-52 p
Acum să urcăm pe scară. Ține-te de \VUgras{balustradă} \textsuperscript{(15)} și
numără treptele. Sunt paisprezece. Iată \VUgras{palierul}
\textsuperscript{(m)}: suntem la \VUgras{etajul} întâi \textsuperscript{(27)}.

%%K t05-tit-9 sub
\VUsaut{3.0mm}
\VUpk{10.2pt}{40}{Etajul întâi.}
\VUsaut{3.0mm}

%%K t05-01-53 p
În fața scării --- \VUgras{closetul} \textsuperscript{(20)} și \VUgras{camera de
toaletă} \textsuperscript{(j)}. La dreapta o \VUgras{cameră de culcare}
frumoasă \textsuperscript{(g)} cu două ferestre spre \VUgras{fațada}
\textsuperscript{(1)} casei. La stânga --- \VUgras{baia} \textsuperscript{(1)} și
\VUgras{camera de oaspeți} \textsuperscript{(i)} cu o \VUgras{alcovă} mare
\textsuperscript{(h)}. Lângă camera de culcare --- \VUgras{camera copiilor}
\textsuperscript{(k)}. Dă spre o \VUgras{terasă} largă \textsuperscript{(21)}. Sub
această terasă un alt closet \textsuperscript{(20)}, iar lângă zid ---
\VUgras{clopotul} \textsuperscript{(22)}, care sună la masă.

%%K t05-01-54 p
\textsc{Ioan}. --- Să urcăm la \VUgras{etajul al doilea}.

%%K t05-01-55 p
\textsc{Arhitectul}. --- Etaj al doilea nu este. Sub acoperiș sunt numai
\VUgras{mansardele} \textsuperscript{(33)} și podul \textsuperscript{(34)}.
Ocolul s-a sfârșit, putem coborî.

%%K t05-01-56 p
\textsc{Ioan}. --- Mulțumesc, domnule, a fost foarte interesant. Îi voi povesti
tatei tot ce am văzut.
\VUsaut{4.0mm}
\VUfilet{22mm}
